\documentclass{article}

%%% Package List

\usepackage{datetime} % Dates.
\usepackage{fancyhdr} % Headers.
\usepackage{lastpage} % Page References.

\usepackage{graphicx} % Inserting Images.
\usepackage{float} % Postitioning Images.

\usepackage{dirtytalk} % Speech & Quote Marks.
\usepackage{hyperref} % Hyperlinks and References.
\usepackage{xcolor} % Coloured Text.
\usepackage{enumitem} % Letter Bullets. 

\usepackage{amssymb} % Maths Symbols.
\usepackage{amsmath} % Maths Tools.

\usepackage{algpseudocode} % Pseudocode.
\usepackage{algorithm} % Algorithms.
\usepackage{minted} % Code Snippets.

\usepackage[margin=1in]{geometry} % Reduces the Page Margins.

%%%



%%% Custom Formats and Commands

\newcommand{\templateversion}{1.8.0} % Do not edit this variable.

\hypersetup{colorlinks=true, urlcolor=black, linkcolor=black}


% Date Formats.
\newdateformat{monthyeardate}
{%
    \monthname[\THEMONTH] \THEYEAR
}

\newdateformat{yeardate}
{%
    \THEYEAR
}


% Colouring and Spacing.
\newcommand{\colour}[2] {\color{#1}#2\color{black}}
\newcommand{\separate}[1] {\vspace{#1 pt}\noindent}


% Simplified Maths Notation.
\newcommand{\R}{\mathbb{R}}
\newcommand{\N}{\mathbb{N}}
\newcommand{\C}{\mathbb{C}}
\newcommand{\Q}{\mathbb{Q}}
\newcommand{\Z}{\mathbb{Z}}
\newcommand{\I}{\mathbb{I}}

\newcommand{\hessian}[1]{\mathbf{H}_{#1}}
\newcommand{\jacobian}[1]{\mathbf{J}_{#1}}
\newcommand{\lagrange}{\mathcal{L}}


% Custom Formats.
\newcommand{\definition}[2]{\textbf{Definition #1}: #2\label{def:#1}}
\newcommand{\example}[2]{\textbf{Example #1}: #2\label{ex:#1}}
\newcommand{\conjecture}[2]{\textbf{Conjecture #1}: #2\label{cnj:#1}}
\newcommand{\numberedentry}[3]{\textbf{#2 #1}: #3\label{entry:#1}}
\newcommand{\proof}[2]{\textit{Proof}: #2 \begin{flushright}\(\square\)\end{flushright}\label{prf:#1}}
\newcommand{\theorem}[3]{\textbf{Theorem #1}: #2 \par\separate{3}\proof{thm:#1}{#3}\label{thm:#1}}
\newcommand{\lemma}[3]{\textbf{Lemma #1}: #2 \par\separate{3}\proof{lem:#1}{#3}\label{lem:#1}}
\newcommand{\proposition}[3]{\textbf{Proposition #1}: #2 \par\separate{3}\proof{prop:#1}{#3}\label{prop:#1}}
\newcommand{\standardproblem}[3]{\textbf{Standard Problem #1}: #2 \par\separate{3}\textit{Solution}: #3\label{prob:#1}}


% Join Symbols.
\newcommand{\ojoin}{\setbox0=\hbox{$\bowtie$}%
    \rule[-.02ex]{.25em}{.4pt}\llap{\rule[\ht0]{.25em}{.4pt}}}
\newcommand{\insertinnerjoin}{\mathbin{\bowtie\mkern-5.8mu}\;}
\newcommand{\insertleftouterjoin}{\mathbin{\ojoin\mkern-5.8mu\bowtie}}
\newcommand{\insertrightouterjoin}{\mathbin{\bowtie\mkern-5.8mu\ojoin}}
\newcommand{\insertfullouterjoin}{\mathbin{\ojoin\mkern-5.8mu\bowtie\mkern-5.8mu\ojoin}}


% Advanced Formats.
\newcommand{\insertimage}[3]
{
    \begin{figure}[H]
        \centering
        \includegraphics[width=#3\linewidth]{#1}
        \caption{#2}
        \label{fig:#2}
    \end{figure}
}

\newcommand{\insertalgorithm}[4]
{
    \begin{algorithm}[H]
        \caption{#1}
        \begin{algorithmic}
            \State \textbf{Input}: #2
            \State \textbf{Output}: #3
            \separate{7}#4
        \end{algorithmic}
        \label{alg:#1}
    \end{algorithm}
}

\newcommand{\inserttable}[4]
{
    \begin{table}[H]
        \centering
        \begin{tabular}{|#1|}
            \hline #3 \\ \hline #4 \\ \hline
        \end{tabular}
        \caption{#2}
        \label{tab:#2}
    \end{table}
}

%%%



%%% Document Setup

% Title Page.
\title{\thetitle\\[0.5em]\large{\thesubtitle}}
\author{\theauthor}
\date{\thedate} 


% Custom Page Header.
\fancyhead{}
\fancyhead[R]{\thetitle\text{ - }\thesubtitle}

% Custom Page Footer.
\fancyfoot{}
\fancyfoot[L]{Written by \theauthor \\ 
\textit{\LaTeX \,Template v.\templateversion: Copyright \date{\yeardate\today}, Matthew Emerson}} % Do not edit.
\fancyfoot[R]{Page \thepage \, of\, \pageref{LastPage}}

%%%



%%% Page Variables

\newcommand{\thetitle}{Discrete Mathematics} % Insert the document title here.
\newcommand{\thesubtitle}{MATH1031 - Durham University} % Insert the document subtitle here.
\newcommand{\theauthor}{Matthew Emerson} % Insert the name(s) of the author(s) here.
\newcommand{\thedate}{\monthyeardate\today} % Replace this if you want a static date.

%%%



%%% Setting Up Pages.

\pagestyle{empty}\begin{document}
\maketitle\thispagestyle{empty} % Title Page.
\setcounter{tocdepth}{2}
\newpage\tableofcontents % Contents Page.

\clearpage\setcounter{page}{1}\newpage % Excludes the title and contents page from page count.
\pagestyle{fancy}

%%%



%%% Document Content

\section{Introduction to Discrete Mathematics}

\subsection{Arrangements and Permutations}

\subsubsection{Arrangements}

\definition{1.1.1.1}{An \underline{arrangement} is a sequence of items in which the order of the items in the sequence matters.}

\separate{7}\definition{1.1.1.2}{The \underline{length} of an arrangement is the number of items in the sequence.}

\separate{7}\standardproblem{1}{How many arrangements of length \(k\) can be made from elements of a set \(S\) of size \(|S|=n\)?}
{\(n^k\) as arrangement has \(k\) places, each of which can be occupied by any one of \(n\) entries.}

\separate{7}\example{1.1.1.1}{How many four-letter \say{words} can be made from the standard A-Z alphabet?

\separate{3}\textit{Solution}: \(26^4 = 456,796\) words.}

\separate{7}\standardproblem{2}{How many arrangements of length \(k\) with no repeats can be made from elements of a set \(S\) of size \(|S|=n\)?}
{\begin{equation} 
    n(n-1)(n-2)...(n-k+1) = \frac{n!}{(n-k)!},\qquad k \leq n    
\end{equation}
As we look at the available choices for each successive item in the list, the number of options is reduced by one each time.}

\separate{7}\definition{1.1.1.3}{The \underline{multiplication principle} states that, for an arrangement of length \(k\), if for \(i \in \{1,\;2,\;...,\;k\}\) there are \(n_i\) choices for the \(n^{th}\) item, regardless of previous choices, then the number of possible arrangements is \(n_1n_2...n_k\).}

\subsubsection{Permutations}

\definition{1.1.2.1}{A \underline{permutation} on a finite set \(S\) is an arrangement of length \(|S|\) that uses each item from \(S\) exactly once.}

\separate{7}\proposition{1.1.2.1}{If \(|S|=n\), then there are \(n!\) permutations on \(S\).}
{This is \hyperref[prob:2]{Standard Problem 2} with \(n=k\).}

\separate{7}\definition{1.1.2.2}{We generalise this concept as, for a non-negative integer \(r\), an \underline{\(r\)-permutation} on \(S\) is an arrangement of \(r\) distinct elements from \(S\). We write \(P(n,r)\) for the number of \(r\)-permutations on a set of size \(n\), which, by \hyperref[prob:2]{Standard Problem 2}, we have,
\begin{equation}
    P(n,r) = n(n-1)...(n-r+1) = \frac{n!}{(n-r)!},\qquad 0 \leq r \leq n
\end{equation}}

\subsubsection{Counting Examples}

\definition{1.1.3.1}{The \underline{addition principle} states that, for two disjoint, finite sets \(A\) and \(B\). Then,
\begin{equation}
    A \cap B = \emptyset \qquad |A \cup B| = |A| + |B|
\end{equation}
More generally, if \(A_1,\;A_2,\;...,\;A_n\) are pairwise disjoint, then
\begin{equation}
    |A_1 \cup A_2 \cup ... \cup A_n| = |A_1| + |A_2| + ... + |A_n|
\end{equation}}

\subsection{Combinations}

\subsection{Arrangements and Combinations with Repetitions}

\section{Combinatorics}

\subsection{The Three Principles}

\subsection{Recurrence Relations}

\subsection{Generating Functions}

\section{Graph Theory}

\subsection{What is a Graph?}

\subsection{Eulerian Graphs}

\subsection{Planar Graphs}

\end{document}

%%%
